% SHARD-ID: CA-52-QUBIT-CANDIDATE-SELECTION-ROADMAP
% SHARD-TITLE: Candidate Hamiltonian Selection Roadmap
% SHARD-SUMMARY: Explains how algebraic gates and SDP levels prioritize real qubit Hamiltonian candidates for later scaling-limit work.
% SHARD-SUMMARY: Feasibility is treated as survival of a finite exclusion test, never as evidence of symmetry.
% SHARD-SUMMARY: Lists the next implementation steps: witness scans, run-bundle format, and physical source registration for named model families.
% SHARD-KEYWORDS: candidate selection, qubit Hamiltonian, screening, SDP hierarchy, Poincare, Virasoro

\section{Candidate Hamiltonian Selection Roadmap}
\label{sec:qubit-candidate-selection-roadmap}

\subsection{Status, Sources, and Dependencies}

\textbf{Status: roadmap.}  This shard connects the implementation to the north
star: identifying Hamiltonians worth expensive continuum analysis.

Dependencies: CA-15, CA-42--CA-51.

% Source: references/lattice-symmetry/KooSaleur1993/source/9312156.tex:1030--1108
%   -- finite lattice Virasoro non-closure warning.
% Source: references/text/CFTFromLatticeFermions.txt:376--388 -- vacuum /
%   ground-state subtraction in Koo-Saleur scaling.

\subsection{Screening Order}

For a fixed nearest-neighbour qubit density \(h_{\alpha\beta}\), use:

\begin{enumerate}
\item current collapse \(p=0\);
\item conservation witness \(A=Du\);
\item boost witness \(B-u-\lambda h+\mu I=Dw\) with \(\lambda>0\);
\item finite moment SDP feasibility for the chosen residuals;
\item only then, low-energy/scaling and Koo-Saleur/Virasoro tests.
\end{enumerate}

The point is to reject impossible candidates cheaply.  A Hamiltonian surviving
all finite checks is not accepted as symmetric; it is merely promoted to the
next, more expensive analysis.

\subsection{Next Implementation Steps}

The immediate next step is a run format for candidate scans: input
\(h_{\alpha\beta}\), current status, witness data, level data, solver status,
and a headline verdict.  Continuous families require external parameter scans
or a later polynomial hierarchy.  Named physical model families need registered
sources before the report can claim criticality, universality class, or
continuum CFT content.
